
\uline{d'%
autre part} (IV 2.11) \uline{on a un homomorphisme}
@@M0@@
\uline{Ceci posé}, \uline{le composé}
@@M1@@
\uline{est nul}.

En effet, @@M2@@ est une immersion régulière de codimension @@M3@@, dont le faisceau conormal @@M4@@ est \emph{libre} de rang @@M5@@, car @@M6@@ peut se décrire par une équation globale @@M7@@ (où @@M8@@ est une uniformisante de @@M9@@). On a donc @@M10@@, et on conclut grâce à 7.1.

Nous supposerons dans toute la suite @@M11@@ \emph{noethérien} et \emph{plat sur} @@M12@@. Appliquons maintenant la suite exacte d'homotopie (IX 1.1) à la situation : @@M13@@, muni de l'ouvert @@M14@@ et du sous-schéma fermé @@M15@@ complémentaire :
@@M16@@
où @@M17@@ est l'inclusion, qui est une immersion ouverte. Utilisant 7.2, on trouve :

\textbf{Proposition 7.3.} \uline{Avec les notations de} 7.2, @@M18@@ \uline{étant supposé noethérien}, \uline{il existe un unique homomorphisme de groupes}
@@M19@@
(\uline{dit} «~homomorphisme de spécialisation~») \uline{rendant commutatif le diagramme}
@@M20@@



% Local macro additions for Exposé X continuation
\providecommand{\Kalg}{K_{\mathrm{alg}}}
\providecommand{\Gralg}{\operatorname{Gr}_{\mathrm{alg}}}
\providecommand{\Grnum}{\operatorname{Gr}_{\mathrm{num}}}
\providecommand{\Knum}{K_{\mathrm{num}}}
\providecommand{\Cl}{\operatorname{Cl}}
\providecommand{\Ch}{\operatorname{Ch}}
\providecommand{\Nm}{\operatorname{N}}
\providecommand{\colim}{\varinjlim}
\providecommand{\Hens}{\operatorname{Hens}}
\providecommand{\Spec}{\operatorname{Spec}}
\providecommand{\coker}{\operatorname{coker}}


